\section{物权的基本原理}
\begin{enumerate}
  \item 4 大内容。\zd{所有权}、用益物权、\zd{担保物权}、准物权（占有）。
  \item \zd{标的物标的物标的物}。
  \item 分类。
        \begin{enumerate}
          \item \zd{动产和不动产}。
          \item 原物和孳息。派生关系。产出 + 分离。主要强调在担保制度、所有权变动中规则的应用。
                \begin{enumerate}
                  \item 天然孳息。\al{下的蛋或崽、掉落的果实}。
                  \item 法定孳息。\al{利息、租金}。
                  \item 射性孳息。\al{彩票、抽奖奖品、保险金}。
                  \item 孳息。\zd{依原物的通常利用方式}从该物中产生的收益。\al{投入股市的收益 \cuo}。
                \end{enumerate}
          \item 主物和从物。辅助关系。
                \begin{enumerate}
                  \item 辅助主物的使用。\al{房屋和车位、空调和遥控器}。
                  \item 非主物的构成部分。\al{汽车和轮胎、汽车和备胎、房屋和房门}。
                  \item \zd{从物与主物要求是同一个人}。\al{甲的手机和乙的充电器}。
                \end{enumerate}
        \end{enumerate}
  \item 3 大效力。\zd{追及、排他、优先}。
        \begin{enumerate}
          \item 追及。\al{甲表辗转于乙、丙、丁}。
          \item 排他。一物一权。
          \item \zd{优先于债权受偿}。\al{甲借乙 1W，以表做担保借丙 1W，到期均无法偿还借款}。
        \end{enumerate}
  \item 物权的保护。请求权基础不同，可以单独适用，可以并用。\al{甲抢走乙的表，不慎摔坏}。
        \begin{enumerate}
          \item 物权请求权（救济权）= 物权（原权利）+ 请求权。\zd{依附于物权的独立请求权}。
                \begin{enumerate}
                  \item \zd{返还原物请求权。所有权 + 无权占有}。考的最多。
                        \al{甲卖房给乙，乙支付了价款，甲交付了房屋但未登记。丁做了简单装修将房屋出租。甲无权请求乙返还房屋}。
                  \item 排除妨害请求权。
                  \item 消除危险请求权。物权\zd{受到妨害的现实可能性}。\al{邻居墙摇摇摇欲坠}。
                  \item 恢复原状请求权。
                \end{enumerate}
          \item 债权请求权。权利人\zd{基于债的关系而产生}的、请求特定人为特定行为的权利。
        \end{enumerate}
\end{enumerate}
